\begin{abstract}
随着城市社区老龄化程度持续加深，嵌入式社区养老服务站建设需要同时兼顾服务可达性、容量可承接性、老人满意度和机构可持续运营。针对某成熟街道 10 个连片小区的养老服务站建设与优化问题，本文围绕需求预测、选址配置、补贴定价和鲁棒评价建立分层优化模型。

针对问题一，本文将老人划分为自理、半失能和失能三类健康状态，构建考虑自然死亡、新增老人和失能不可逆转移的多状态 Markov 递推模型。进一步建立服务需求矩阵，将第 5 年末老人结构映射为助餐、日间照料、上门护理、康复理疗、助浴和紧急救助六类服务需求，并引入消费能力约束修正理论需求。计算结果表明，第 5 年末全区老人总量约为 7578 人，其中自理老人约 4982 人、半失能老人约 1471 人、失能老人约 1126 人；消费约束后六类服务有效月需求分别为助餐 118459 次、日间照料 76359 次、上门护理 19614 次、康复理疗 21185 次、助浴 6538 次和紧急救助 4944 次。

针对问题二，本文建立感知选择与后台容量调度相结合的服务站选址配置模型。模型以空间覆盖率为第一目标，以容量满足率、有效服务率和服务加权满意度为后续目标，综合考虑 120 万元建设预算、1000 米有效服务半径、小型、中型、大型服务站容量约束和满意度评分规则。基准情景下得到的前沿方案为在 C、E、G、I、J 小区建设服务站，规模分别为小型、小型、小型、小型和大型，总建设成本为 117 万元，空间覆盖率达到 100\%，容量满足率为 0.8499，有效服务率为 0.7439，服务加权满意度为 0.8753。

针对问题三，本文在既定站点布局下建立补贴约束定价优化模型。根据价格满意度的分段规则，将连续定价问题转化为有限临界价格组合枚举问题；同时纳入政府按有效服务人次补贴、单站日补贴上限、紧急救助公益免费和机构利润率不超过 8\% 等约束。基准情景下，年度政府补贴总额为 240.90 万元，补贴后年度利润合计为 33.82 万元，平均利润率为 7.2691\%，最低利润率为 5.9786\%，所有服务站均满足保本微利约束，价格满意度为 1.0000。

针对问题四，本文设置老人增长、失能转移、固定成本和建设预算四类扰动情景，采用同一模型重新求解选址和定价方案。结果表明，老人增长、转移概率上升和固定成本上升情景下，最优站点方案均保持为 C小、E小、G小、I小、J大；当建设预算提高至 140 万元时，最优方案调整为 A小、B大、G大、I中，容量满足率提升至 1.0000，有效服务率提升至 0.9065。

本文的主要特点在于：一是将潜在服务需求与消费约束后有效需求分离，避免容量配置虚高；二是构造覆盖、容量、满意度和利润一体化的服务站选址评价框架；三是利用临界点枚举处理分段价格满意度，保持求解过程可复现；四是以情景重求解而非单纯扰动讨论检验模型稳健性。

\keywords{嵌入式养老服务站\quad Markov 递推\quad 容量约束选址\quad 满意度评价\quad 补贴定价\quad 鲁棒优化}
\end{abstract}
